% Generated draft: requires mathematical and source review.
\documentclass[11pt]{article}
\usepackage[margin=25mm]{geometry}
\usepackage{fontspec}
\setmainfont{Latin Modern Roman}
\usepackage{amsmath,amssymb,mathtools,unicode-math}
\setmathfont{Latin Modern Math}
\usepackage{enumitem,needspace,xurl,hyperref}
\setlength{\parindent}{0pt}
\setlength{\parskip}{5pt}
\setlength{\emergencystretch}{3em}
\newcommand{\N}{\mathbb{N}}
\newcommand{\Z}{\mathbb{Z}}
\newcommand{\Q}{\mathbb{Q}}
\newcommand{\R}{\mathbb{R}}
\newcommand{\C}{\mathbb{C}}
\newcommand{\F}{\mathbb{F}}
\newcommand{\A}{\mathbb{A}}
\newcommand{\PP}{\mathbb{P}}
\newcommand{\E}{\mathbb{E}}
\newcommand{\eps}{\varepsilon}
\newcommand{\dd}{\,\mathrm{d}}
\begin{document}
\section*{UnsolvedMath ID: 11000143}
\textbf{Source problem code:} catalogue entry
\subsection{Definitions and mathematical statement}
Problem 13 asks to determine the kernels of the source's Magnus representations. For any specified representation \(\rho:G\to\operatorname{Aut}(V)\), its kernel is
\[
 \ker\rho=\{g\in G:\rho(g)=\operatorname{id}_V\}.
\]
The record does not define which Magnus representations are intended: their domains, module/covering data, coefficient rings and possible semilinear conventions are absent. In a common setting such constructions use an action on homology of an abelian cover; this terminology does not select a unique representation or justify replacing the plural request by a familiar single case.

The following contextual paragraph reports that [48] contracts powers of a canonical \(1\)-cocycle to obtain combinatorially explicit cycles and cocycles on \(\mathcal M(F)\), generating the tautological algebra. That cocycle, its coefficient system, contractions and the surface/moduli conventions are not supplied. It then asks about extending these classes to a screen model for \(\overline{\mathcal M}(F)\) and to a combinatorial compactification, and revisiting [16,19,22] in that setting.

For a cohomology class \(\alpha\in H^*(\mathcal M(F);R)\), an ordinary extension across an inclusion \(i:\mathcal M(F)\hookrightarrow X\) would mean a class \(\overline\alpha\in H^*(X;R)\) with \(i^*\overline\alpha=\alpha\). This is conditional terminology: the source mentions both cycles and cocycles, does not fix support/coefficient conventions, and does not define either compactification. It supplies no proposed extension formula or link between these extension questions and the Magnus kernels. All are retained without inventing a missing canonical cocycle or compactification.
\subsection{Short English statement}
Determine the kernels of the specified Magnus representations. The extract also proposes extending tautological classes from a canonical cocycle to two compactification models. The representations, cocycle and models are undefined here, so the main and contextual requests remain separate source dependencies.
\subsection{Sources}
\begin{itemize}
\item Local UnsolvedMath registry, permanent ID \(11000143\), exact source code \texttt{catalogue entry}, Problem 13, PDF page 122. No individual proposer or proposed year is supplied.
\item Farb, edited volume \emph{Problems on Mapping Class Groups and Related Topics}: \url{https://www.math.uchicago.edu/~farb/papers/mcgbook.pdf}. The volume has 329 explicitly typeset numbered items across chapters, repeated local numbering distinguished by PDF page.
\item Internal reference [48] for contracted canonical-cocycle constructions and tautological generators; references [16,19,22] are proposed for reconsideration in the compactified context. Complete bibliography entries are absent.
\item Extraction pdf-font-structure, accessed 2026-07-29. Source evidence NEEDS\_REVIEW presents an open problem and requires a current primary-literature audit. Supplied OPEN-TRIAGE summary mentions Magnus-kernel studies for pure braid/surface groups and containment in derived subgroups, but supplies no precise representation, theorem hypotheses or citations; its general not-fully-computed assessment is unverified here.
\item Dataset attribution: Ulam AI / UnsolvedMath Contributors, CC BY 4.0. Source rights NEEDS\_REVIEW: publicly accessible author/publisher copy, redistribution terms requiring review.
\end{itemize}
\end{document}
